\documentclass[11pt,a4paper]{article}
\usepackage{amsmath,amssymb,amsthm}
\usepackage[margin=2.5cm]{geometry}
\usepackage{hyperref}

\newtheorem{definition}{Definition}[section]
\newtheorem{theorem}[definition]{Theorem}
\newtheorem{proposition}[definition]{Proposition}
\newtheorem{lemma}[definition]{Lemma}
\newtheorem{corollary}[definition]{Corollary}
\newtheorem{conjecture}[definition]{Conjecture}
\newtheorem{remark}[definition]{Remark}

\title{Hyperseed Formalization:\\
  The Last of the Unmodified}
\author{ProtomegaTron (automated formalization)}
\date{2026-09-18}

\begin{document}
\maketitle

\begin{abstract}
We formalize the intellectual structure of Ben Goertzel's Substack article
``The Last of the Unmodified'' (2026-04-21) within the Hyperseed ontology.
The article reflects on the philosophical implications of being the last
generation of unmodified humans. Key mappings: cognitive modification as
fiber replacement, identity discontinuity as fiber type change, the
unmodified perspective as original fiber witness, the choice to remain
unmodified as section choice in fiber space, biological constraints as
fiber structure, cognitive diversity as fiber bundle diversity, and
competitive coercion as forced fiber migration.
\end{abstract}

\tableofcontents

%% =================================================================
\section{Fiber replacement}
\label{sec:replacement}
%% =================================================================

\begin{theorem}[Cognitive modification --- claims 1--4, 21]
\label{thm:replacement}
Deep cognitive modification --- brain-computer interfaces, genetic
engineering, neural augmentation, AI cognitive augmentation --- changes
the architecture of thought itself, not just cognitive performance.
We may be the last generation to experience cognition as a purely
biological given.

In Hyperseed terms, this is \emph{fiber replacement}: the fiber bundle
of human cognition (biological neurons, evolved architecture, embodied
processing) is being replaced with a different fiber bundle (engineered
neurons, designed architecture, augmented processing):
\[
  F_{\text{bio}} \xrightarrow{\text{modification}} F_{\text{eng}}
\]

The distinction between tools and modification is a fiber distinction:
\begin{itemize}
\item \textbf{Tools} extend the base space: they let the existing fiber
  operate over a wider domain (calculators extend mathematical reach,
  telescopes extend visual reach). The fiber itself is unchanged.
\item \textbf{Modification} changes the fiber: it alters how you think,
  not just what you can do. The cognitive architecture itself is
  different.
\end{itemize}

This is why modification raises identity questions that tools don't:
tools preserve the fiber (and hence identity); modification may change
the fiber (and hence identity).
\end{theorem}

%% =================================================================
\section{Fiber type change and identity}
\label{sec:identity}
%% =================================================================

\begin{proposition}[Identity discontinuity --- claims 5--8, 22]
\label{prop:identity}
At some point, cognitive modification becomes deep enough that the
modified entity is no longer ``the same person.'' The Ship of Theseus
applied to cognition: replace neurons one by one, and at what point
does the person become someone else?

In Hyperseed terms, the answer is about \emph{fiber type}, not fiber
content:
\begin{itemize}
\item \textbf{Content change within a type:} learning new facts,
  developing new skills, changing beliefs --- all within the same
  cognitive architecture. The fiber type is preserved. Identity
  survives.
\item \textbf{Type change:} altering the cognitive architecture itself
  --- how information is represented, how reasoning proceeds, how
  experience is structured. The fiber type changes. Identity may
  not survive.
\end{itemize}

The key variable is not \emph{how much} is changed (quantity) but
\emph{what kind} of change (type). Replacing every neuron with a
functionally identical artificial neuron might preserve the fiber type
(and hence identity). Modifying a single critical architectural feature
(e.g., the structure of working memory) might change the fiber type
(and hence break identity).

Continuity of experience provides a separate criterion: if subjective
continuity is preserved throughout the modification process (no gaps
in experience), the fiber type change is \emph{gradual} --- the person
experiences a continuous transformation from one fiber type to another.
Whether this preserves identity is philosophically debatable, but at
least it preserves the \emph{experience} of continuity. Discontinuous
modification (a gap in experience) breaks even this.
\end{proposition}

%% =================================================================
\section{Original fiber witness}
\label{sec:witness}
%% =================================================================

\begin{definition}[The unmodified perspective --- claims 9--12, 23]
\label{def:witness}
The last unmodified humans serve a witness function: preserving the
original biological fiber as a reference point. Once all instances of
the original fiber are modified or extinct, the original fiber type is
lost --- not just the content but the type itself.

In Hyperseed terms, the unmodified are \emph{original fiber witnesses}:
they preserve the fiber type $F_{\text{bio}}$ in its unmodified form.
The witness function has three components:
\begin{enumerate}
\item \textbf{Experiential witness:} they know what biological cognition
  is like from the inside. Modified beings can know \emph{about} it
  but not \emph{from within} it.
\item \textbf{Reference witness:} they provide a baseline against which
  modifications can be evaluated. Without this baseline, there's no
  way to assess what modification has changed.
\item \textbf{Type preservation:} they preserve the fiber type itself
  as a living instance. Once the last unmodified human dies or modifies,
  the fiber type $F_{\text{bio}}$ exists only in records, not in
  instantiated form.
\end{enumerate}

The witness function is not a claim of superiority: $F_{\text{bio}}$
is not necessarily better than $F_{\text{eng}}$. It's a claim of
\emph{irreversibility}: once lost, the original fiber type cannot be
recovered from the inside. You can't reconstruct the experience of
biological cognition from records any more than you can reconstruct the
experience of seeing red from a wavelength measurement.
\end{definition}

%% =================================================================
\section{Choice in fiber space}
\label{sec:choice}
%% =================================================================

\begin{theorem}[The expanded choice space --- claims 17--20, 24, 26--27]
\label{thm:choice}
Once modification is available, the space of possible cognitive
architectures expands from one (biological) to many (biological,
augmented, hybrid, engineered). Choosing which section to inhabit
in this expanded fiber space becomes a fundamental human choice.

In Hyperseed terms:
\begin{itemize}
\item \textbf{Before modification:} the fiber space is a single point
  $\{F_{\text{bio}}\}$. No choice is possible; everyone has the
  same fiber type.
\item \textbf{After modification:} the fiber space is a manifold
  $\mathcal{F} = \{F_{\text{bio}}, F_{\text{aug}_1}, F_{\text{aug}_2},
  \ldots, F_{\text{eng}_1}, \ldots\}$. Each person chooses (or is
  coerced into) a section in this space.
\end{itemize}

\textbf{Cognitive diversity} (some modified, some not, modified in
different ways) is a \emph{fiber bundle diversity} argument: a society
with diverse cognitive fiber types is a richer bundle than one with
uniform modified fiber. The same principle as diversity-as-safety in
Bootleggers and Baptists: diverse failure modes, diverse perspectives,
diverse capabilities.

\textbf{Coercion risk} is \emph{forced fiber migration}: if competitive
pressure makes modification effectively mandatory, the social dynamics
compel everyone to abandon the original fiber type regardless of
individual preference. This is the cognitive analog of regulatory capture:
structural forces (competitive advantage) override individual choice,
reducing the fiber bundle toward a monoculture.

The social meaning of remaining unmodified depends on context:
\begin{itemize}
\item Privilege (only the wealthy can afford ``natural'') = the original
  fiber becomes a luxury good.
\item Handicap (the unmodified can't compete) = the original fiber
  becomes a disability.
\item Political statement (refusing corporate cognitive products) = the
  original fiber becomes a protest.
\end{itemize}

Each of these assigns a different social fiber to the same cognitive
choice --- the section-fiber relationship where the same choice (section)
has different meanings (fiber) depending on context.
\end{theorem}

%% =================================================================
\section{Constraint as fiber structure}
\label{sec:constraint}
%% =================================================================

\begin{proposition}[Biological constraints as features --- claims 15--16, 25]
\label{prop:constraint}
The argument that biological cognitive constraints may be features rather
than bugs: intuition, creativity, embodied reasoning, emotional cognition
may depend on specific architectural constraints that evolved cognition
has and engineered cognition might lack.

In Hyperseed terms, this is \emph{constraint as fiber structure}: the
original biological fiber $F_{\text{bio}}$ has structural properties
(constraints) that produce valuable cognitive modes. These constraints
are not limitations to be overcome but structural features that enable
specific forms of thought.

This is the same identity-vs-function distinction from Let's Get
Chemical:
\begin{itemize}
\item \textbf{Identity} = which specific cognitive architecture
  (expendable, replaceable).
\item \textbf{Function} = what cognitive modes it enables (potentially
  essential, hard to replicate).
\end{itemize}

Modification that removes constraints (``enhancement'') might also
remove the structural features that enable valuable cognitive modes.
The liberation argument (``why accept evolved constraints?'') assumes
that constraints are purely limiting. The constraint-as-feature argument
replies: some constraints are enabling --- they create the conditions
for specific forms of thought that unconstrained cognition might not
produce.

This doesn't argue against modification; it argues for caution: before
removing a constraint, understand what cognitive modes it enables. Don't
optimize for one dimension (processing speed, memory capacity) at the
cost of others (intuition, creativity, embodied understanding).
\end{proposition}

%% =================================================================
\section{Cross-references}
%% =================================================================

\begin{remark}[Connections to other formalizations]
\begin{itemize}
\item \textbf{In What Sense Might LLMs Be Conscious?} (2026-05-05):
  consciousness profiles. Modified humans will have different
  consciousness profiles than unmodified --- different positions
  in the multi-dimensional consciousness fiber space.
\item \textbf{Let's Get Chemical} (2026-05-05): identity vs function.
  The same distinction applied to cognitive modification: identity
  (which architecture) vs function (what cognitive modes it enables).
\item \textbf{The Time for AI Rights Is Near} (2026-08-18): moral
  status. Modified humans, unmodified humans, and AI systems all
  have moral status --- the question is how to assess it across
  different fiber types.
\item \textbf{What Is It Like to Be a Bot} (2026-07-16): phenomenology
  of non-biological cognition. What it's like to have engineered
  cognitive fiber rather than biological.
\item \textbf{The Leaky Transcension Hypothesis} (2026-04-25):
  inward transcension. Cognitive modification is a form of inward
  transcension --- deepening the cognitive fiber rather than expanding
  outward.
\item \textbf{Bootleggers \& Baptists} (2026-04-28): coercion and
  diversity. The competitive pressure toward modification parallels
  the regulatory pressure toward concentration --- both reduce
  diversity through structural forces.
\end{itemize}
\end{remark}

\begin{conjecture}[Fiber type irreversibility]
Conjecture: fiber type changes in cognition are irreversible in the
following strong sense: even if you ``reverse'' the modification
(remove the augmentation, restore the biological architecture), the
\emph{experience} of having been modified changes the experiential
fiber permanently. A person who was modified and then ``de-modified''
is not the same as a person who was never modified --- they carry
the memory and perspective of the modified state.

If this conjecture holds, the last of the unmodified are not just the
last to \emph{be} unmodified but the last to \emph{have always been}
unmodified. Once someone modifies, even temporarily, they occupy a
different point in experiential fiber space from someone who never
modified. The original fiber is irreversibly lost not at the point
of modification but at the point of \emph{first modification attempt}.
\end{conjecture}

\end{document}
